\section{Conclusion}
\label{sec:conclusion}

We presented \textbf{PRISM}, a multimodal sentiment model that compresses
each modality with variational information bottlenecks, routes a per-sample
primary stream with a Gumbel--Softmax selector, and fuses auxiliaries through
IB-guided cross-attention and adaptive gates---committing the final score to
the \emph{enhanced primary} bottlenecks only, in the same spirit as
primary-modality-centric designs in MODS \citep{mods2026align} and IB analyses
in CyIN \citep{cyin2025align}, but with a concrete, code-matched
implementation.
IB confidence provides a direct, learnable handle on which key positions
auxiliary modalities may attend to, while routing supervision in stage~2 keeps
the selector tracking per-modality predictive error rather than freezing to a
static ``language-first'' prior; the design also trades off full variational
deep hierarchies (as might be inspired by ITHP-style fusion
\citep{ithp2024align}) for a lighter diagonal Gaussian and a single fusion
stack, favoring training stability on standard MSA budgets.
On CMU-MOSI, CMU-MOSEI, and a large CH-SIMS v2 Optuna search, PRISM obtains
competitive or improved MAE/accuracy operating points under complete-modality
protocols, and MSA features can reflect demographic and channel bias so that
high-stakes use (e.g., hiring or clinical screening) is \emph{not} supported
by calibration of our objective alone---see \S\ref{sec:limitations}.

\noindent \textbf{Future work}.
We see three directions: (i) token-level reparameterized bottlenecks closer to
CyIN's formulation \citep{cyin2025align} at the MSelector interface;
(ii) missing-modality and OOD evaluation;
(iii) joint scaling with stronger visual/acoustic backbones and fair comparison
to causal IB baselines \citep{jiang2025towards}.
